\section{El teorema de excisión}

% ------------------------------------------------------------------------
\subsection{Enunciado}

\begin{frame}
  \frametitle{El teorema de excisión}

  Para $A\subseteq X$, las cadenas relativas son
  $\Sing_\bullet(X,A)=\Sing_\bullet(X)/\Sing_\bullet(A)$, y la homología
  relativa es $H_n(X,A)=H_n\bigl(\Sing_\bullet(X,A)\bigr)$.

  \begin{theorem}[Excisión]
    Sean $Z\subseteq A\subseteq X$ con $\overline{Z}\subseteq\Int A$.
    Entonces la inclusión de pares
    $(X\setminus Z,A\setminus Z)\hookrightarrow(X,A)$ induce isomorfismos
    \[
      H_n(X\setminus Z,A\setminus Z)\cong H_n(X,A)
    \]
    para todo $n$.
  \end{theorem}

  Es decir, podemos \alert{suprimir} $Z$ sin alterar la homología
  relativa, siempre que $Z$ esté bien dentro de $A$.
\end{frame}

\begin{frame}[allowframebreaks]
  \frametitle{Una formulación equivalente}

  \begin{theorem}[Excisión, segunda forma]
    Sean $A,B\subseteq X$ con $\Int A\cup\Int B=X$. Entonces la inclusión
    de pares $(B,A\cap B)\hookrightarrow(X,A)$ induce isomorfismos
    \[
      H_n(B,A\cap B)\cong H_n(X,A)
    \]
    para todo $n$.
  \end{theorem}

  \framebreak

  \begin{proposition}
    La segunda forma implica la primera.
  \end{proposition}

  \begin{proof}
    Sean $Z\subseteq A\subseteq X$ con $\overline{Z}\subseteq\Int A$, y sea
    $B=X\setminus Z$. Entonces
    \[
      \Int B=X\setminus\overline{Z}\supseteq X\setminus\Int A,
    \]
    luego $\Int A\cup\Int B=X$. Además, $A\cap B=A\setminus Z$, así que la
    inclusión $(B,A\cap B)\hookrightarrow(X,A)$ es la del enunciado de la
    primera forma.
  \end{proof}

  \medskip

  Demostraremos, por tanto, la segunda forma. La hipótesis dice
  exactamente que $\cU=\{A,B\}$ es un recubrimiento de $X$.
\end{frame}

% ------------------------------------------------------------------------
\subsection{Demostración}

\begin{frame}
  \frametitle{Primer paso: un isomorfismo de complejos}

  Sea $\cU=\{A,B\}$, de modo que
  $\CU_\bullet(X)=\Sing_\bullet(A)+\Sing_\bullet(B)$.

  \begin{lemma}
    La inclusión induce un isomorfismo de complejos
    \[
      \Sing_\bullet(B)/\Sing_\bullet(A\cap B)
      \xrightarrow{\ \cong\ }
      \CU_\bullet(X)/\Sing_\bullet(A).
    \]
  \end{lemma}

  \begin{proof}
    Es el segundo teorema de isomorfía,
    $B'/(A'\cap B')\cong(A'+B')/A'$, aplicado en cada grado a
    $A'=\Sing_n(A)$ y $B'=\Sing_n(B)$, teniendo en cuenta que
    $\Sing_n(A)\cap\Sing_n(B)=\Sing_n(A\cap B)$. El isomorfismo es inducido
    por la inclusión, luego conmuta con el borde.
  \end{proof}
\end{frame}

\begin{frame}[allowframebreaks]
  \frametitle{Segundo paso: cadenas pequeñas relativas}

  \begin{lemma}
    La inclusión induce isomorfismos en homología
    \[
      H_n\bigl(\CU_\bullet(X)/\Sing_\bullet(A)\bigr)
      \xrightarrow{\ \cong\ }
      H_n\bigl(\Sing_\bullet(X)/\Sing_\bullet(A)\bigr).
    \]
  \end{lemma}

  \begin{proof}
    Sean $\rho$ y $D$ los del teorema de las cadenas pequeñas. Como
    $\Sing_\bullet(A)\subseteq\CU_\bullet(X)$, se tiene que $\rho$ fija
    $\Sing_\bullet(A)$ y $D$ se anula sobre $\Sing_\bullet(A)$. Por tanto
    $\iota$, $\rho$ y $D$ pasan al cociente por $\Sing_\bullet(A)$:
    \[
      \bar\iota\colon\CU_\bullet(X)/\Sing_\bullet(A)
        \rightleftarrows\Sing_\bullet(X)/\Sing_\bullet(A)\colon\bar\rho,
      \qquad
      \bar D\colon\Sing_n(X,A)\to\Sing_{n+1}(X,A).
    \]

    \framebreak

    Las identidades del teorema de las cadenas pequeñas pasan al cociente:
    \[
      \bar\rho\bar\iota=\id,
      \qquad
      \partial\bar D+\bar D\partial=\id-\bar\iota\bar\rho.
    \]
    Así, $\bar\iota$ es una equivalencia de homotopía de cadenas con inversa
    $\bar\rho$, y por tanto induce isomorfismos en homología.
  \end{proof}
\end{frame}

\begin{frame}
  \frametitle{Conclusión}

  \begin{proof}[Demostración de la excisión]
    La inclusión $\Sing_\bullet(B,A\cap B)\to\Sing_\bullet(X,A)$ inducida por
    $(B,A\cap B)\hookrightarrow(X,A)$ es la composición
    \[
      \Sing_\bullet(B)/\Sing_\bullet(A\cap B)
      \xrightarrow{\ \cong\ }
      \CU_\bullet(X)/\Sing_\bullet(A)
      \xrightarrow{\ \bar\iota\ }
      \Sing_\bullet(X)/\Sing_\bullet(A).
    \]
    El primer morfismo es un isomorfismo de complejos por el primer paso, y
    el segundo induce isomorfismos en homología por el segundo paso. Por
    tanto, la composición induce isomorfismos
    $H_n(B,A\cap B)\cong H_n(X,A)$.
  \end{proof}
\end{frame}

% ------------------------------------------------------------------------
\subsection{Versión reducida}

\begin{frame}
  \frametitle{Excisión en homología reducida}

  Para $A\subseteq X$, la homología relativa reducida
  $\widetilde{H}_n(X,A)$ es la homología de
  $\widetilde{\Sing}_\bullet(X)/\widetilde{\Sing}_\bullet(A)$.

  \begin{corollary}
    En las hipótesis de la excisión, si además $A\setminus Z\neq\varnothing$,
    la inclusión induce isomorfismos
    \[
      \widetilde{H}_n(X\setminus Z,A\setminus Z)\cong\widetilde{H}_n(X,A)
    \]
    para todo $n$.
  \end{corollary}

  \begin{proof}
    Si $A\neq\varnothing$, los complejos aumentados de $X$ y de $A$ coinciden
    en grado $-1$, luego
    $\widetilde{\Sing}_\bullet(X)/\widetilde{\Sing}_\bullet(A)
    =\Sing_\bullet(X)/\Sing_\bullet(A)$ y
    $\widetilde{H}_n(X,A)=H_n(X,A)$. Lo mismo vale para el par
    $(X\setminus Z,A\setminus Z)$, pues $A\setminus Z\neq\varnothing$. Basta
    aplicar el teorema de excisión.
  \end{proof}
\end{frame}
